Human history has been profoundly shaped by waves of pathogen-borne epidemics. They have left enduring marks on our species, ranging from our behaviour and social norms to our cellular and genetic material. A major factor of infectious disease outbreaks is our interactions with other animal species, which have become increasingly frequent with migration, domestication, and ecosystem-eroding technological advancements. Such interactions facilitate spillover events~\cite{vora2022, sikora2025}, in which a pathogen prevalent in a host animal reservoir is transmitted to a nearby human population. Although progress in science, medicine, and public health awareness has greatly mitigated the frequency and impact of \emph{zoonotic} epidemics, the acceleration of globalisation and climate change in the 21st century is likely to change the balance~\cite{carlson2022}. Indeed, recent estimates posit that a COVID-19-style pandemic is likely to occur once every 129 years on average~\cite{marani2021}. Among other factors, ecosystem modifications (e.g., deforestation, loss of biodiversity~\cite{ostfeld2009, keesing2010}) reduce buffers between zoonotic reservoirs and human populations, while global warming facilitates the geographical expansion of disease-carrying vectors such as bats, mosquitoes, and rodents~\cite{li2005, jones2008, ryan2019, johnson2020}. 

Since the early 2000s, global organisations such as the \gls{who} and the World Bank have spearheaded efforts such as the Global Outbreak Alert and Response Network (GOARN) and the Global Pandemic Monitoring Board (GPMB) to strengthen international pandemic coordination~\cite{heymann2001, knobler2004, who2025}. Yet, as illustrated by the recent COVID-19 pandemic, preparedness against such events remains a Herculean task. Beyond logistical and sociopolitical issues, the scientific landscape of \emph{pandemic preparedness} is inherently interdisciplinary~\cite{fefferman2023}. For instance, microbiology is paramount for deciphering the structure and function of pathogenic agents; ecology and veterinary science are essential for studying the dynamics of animal reservoirs; advanced statistics are needed for modelling the progress of epidemics; and information technology is crucial for developing efficient software. Nowadays, \emph{data science} has become a pillar of pandemic preparedness, as research in infectious disease epidemiology increasingly relies on the analysis of large-scale heterogeneous datasets~\cite{leonelli2021}. These include wildlife monitoring footage, pathogenic genomes from wastewater, human mobility patterns, and electronic health records. Moreover, a substantial body of research on the COVID-19 pandemic has highlighted the utility of large-scale, possibly population-wide, datasets to better understand the dynamics of disease transmission, identify risk factors, and inform effective prevention and intervention strategies~\cite{grenfell2004, robishaw2021, vigfusson2021, attwood2022, matteson2023, khurana2024}. Thus, investment in data-driven research is crucial for our preparedness for the next pandemic.

\newpage

This thesis aims to contribute to pandemic preparedness through advances in data science, in which a common underlying theme is the mathematical \emph{vectorisation} of pandemic-related data. Representing complex data in a simple array of numbers is desirable for efficient storage, scalability, and interoperability with various statistical frameworks. More specifically, I first introduce \emph{phylo2vec}~\cite{penn2025}, a method for encoding phylogenetic trees (a representation of evolutionary relationships between entities) into a simple integer-based vector. The primary objective of this method was to enable efficient data compression, rapid tree sampling and to facilitate the use of phylogenetic trees in comparative analyses and potential machine learning applications. The objective of the second study was to explore the potential of protein structures to complement our understanding of viral pathogen evolution. Whereas viral evolution is typically studied through sequence analysis, structures may offer new insight as a more direct reflection of protein function~\cite{crick1958}. Consequently, using betacoronaviruses as an example, I used phylogenetic mixture modelling~\cite{pagel2005} to analyse the evolution of spike glycoproteins, which these viruses use for host cell entry. Lastly, the objective of the third study was to develop a framework for video-based animal tracking~\cite{scheidwasser2025}, to enable early, non-invasive detection of infection in farm settings. To this end, using broiler chickens as an example, I co-developed a feasibility study to detect \textit{Escherichia coli} infections from behavioural features extracted from videography. Together, these contributions highlight the pivotal role of efficient data representation in strengthening our capacity to anticipate, monitor, and respond to future pandemics.
